% ============================================================
% Section 6: Dual Heterophily (Medium length for TNNLS)
% Target: ~1.5 pages
% ============================================================

\section{Two-Hop Recovery as a Heterophily Diagnostic}
\label{sec:heterophily}

Low-homophily datasets exhibit puzzling inconsistency: graph models improve strongly on Chameleon and Squirrel ($h \approx 0.22$), while GCN degrades sharply on Texas and Wisconsin ($h < 0.2$). We test two-hop recovery as one possible explanatory mechanism.

\subsection{The Heterophily Puzzle}

Traditional wisdom suggests GNNs fail on heterophilic graphs because neighbors have different labels. However, this oversimplifies:

\begin{itemize}
    \item \textbf{Texas} ($h = 0.087$): MLP achieves 80.5\%, GCN achieves 56.3\% (GCN harms)
    \item \textbf{Chameleon} ($h = 0.231$): MLP achieves 51.4\%, while the best recorded graph model, H2GCN, achieves 65.5\%
\end{itemize}

Both have similar homophily, yet opposite outcomes. Low homophily alone cannot explain this.

\subsection{2-Hop Recovery Ratio}

\begin{definition}[2-Hop Homophily]
\label{def:2hop}
Let $h_1$ denote standard (1-hop) homophily and $h_2$ denote 2-hop homophily:
\begin{equation}
    h_2 = \mathbb{P}(Y_v = Y_u \mid \text{dist}(u, v) = 2)
\end{equation}
The \textbf{2-hop recovery ratio} is:
\begin{equation}
    R = \frac{h_2}{h_1}
\end{equation}
\end{definition}

\textbf{Intuition}: When $h_1$ is low (direct neighbors differ), high $R$ indicates 2-hop neighbors are likely same-class---information can be recovered via longer paths.

\begin{definition}[Two-Hop Recovery Patterns]
\label{def:types}
Based on the recovery ratio $R$:
\begin{itemize}
    \item \textbf{Strong recovery (legacy Type A)} ($R > 1.5$): 2-hop neighbors are more likely same-class than 1-hop neighbors, indicating a longer-range label pattern.

    \item \textbf{Absent recovery (legacy Type B)} ($R \leq 1.0$): No observed two-hop homophily recovery. This describes the statistic only and does not imply that graph models underperform MLP.
\end{itemize}
\end{definition}

\subsection{Empirical Validation}

\begin{table}[t]
\centering
\caption{Two-Hop Recovery Patterns and Protocol-Specific Best Models}
\label{tab:2hop}
\begin{tabular}{@{}lccccl@{}}
\toprule
Dataset & $h_1$ & $h_2$ & $R$ & Type & Best Model \\
\midrule
% Generated from code/results/dual_heterophily_10seed_candidate.json. Do not edit manually.
Texas & 0.087 & 0.571 & 6.56 & A & H2GCN \\
Wisconsin & 0.192 & 0.425 & 2.21 & A & MLP \\
Cornell & 0.127 & 0.392 & 3.07 & A & H2GCN \\
Chameleon & 0.231 & 0.230 & 1.00 & B & H2GCN \\
Squirrel & 0.222 & 0.197 & 0.88 & B & GCN \\
\bottomrule
%
\end{tabular}
\end{table}

\textbf{Key Finding} (Table~\ref{tab:2hop}): Large recovery ratios identify datasets where two-hop models substantially repair GCN degradation, but they do not guarantee that H2GCN beats a tuned MLP or every competing GNN.

\subsection{H2GCN Performance Analysis}

H2GCN \cite{zhu2020beyond} is designed for heterophilic graphs, using:
\begin{enumerate}
    \item Ego-neighbor separation (preserve original features)
    \item Higher-order neighborhoods (2-hop aggregation)
    \item Combination of intermediate representations
\end{enumerate}

\textbf{Protocol-specific result}:
\begin{itemize}
    \item \textbf{Strong recovery}: H2GCN substantially repairs GCN degradation on Texas, Wisconsin, and Cornell, but exceeds MLP clearly only on Texas in the 10-seed paired analysis.

    \item \textbf{Absent recovery}: H2GCN is best on Chameleon and GCN is best on Squirrel despite recovery ratios at or below one. Thus $R$ is not a sufficient model-selection statistic.
\end{itemize}

\subsection{Why Type B Exists}

Type B heterophily arises when edge formation is \textbf{class-independent}:
\begin{itemize}
    \item \textbf{Actor network}: Co-appearance in films is determined by casting decisions, not actor attributes
    \item \textbf{Wikipedia pages}: Links reflect topic coverage, not page categories
\end{itemize}

In these cases, measured two-hop homophily does not recover the one-hop label disagreement. Other structural signals or architectures remain useful, as shown by H2GCN on Chameleon and GCN on Squirrel in the audited protocol.

\subsection{Practical Diagnostic Checklist}

\begin{tcolorbox}[colback=green!5!white,colframe=green!75!black,title=Heterophilic Dataset Diagnostic]
For datasets with $h < 0.3$:
\begin{enumerate}
    \item Compute 2-hop recovery ratio $R = h_2/h_1$ as a descriptive graph statistic
    \item Use $R>1.5$ to flag strong two-hop recovery, then compare models empirically rather than selecting one from $R$
    \item If $R \leq 1.0$: Treat two-hop recovery as absent and compare MLP with non-two-hop graph baselines
    \item If $1.0 < R < 1.5$: Empirically compare both
\end{enumerate}
\end{tcolorbox}

This checklist is descriptive rather than a selection rule. In the audited five-dataset heterophily subset, the recovery statistic does not identify the single best model.
